% =====================================================================
%  R. Nielsen, LOGIC AND PSYCHOLOGY.  COURSE FOR THE ACADEMIC YEAR 1871-72.
%  Copenhagen: J. H. Schubothe (Louis Klein), [1871].  8 pp.  (Logik og Psychologie. Cursus for
%  Universitetsaaret 1871-72.)  The print holds the section „Psychologie“ only.
%  ENGLISH TRANSLATION, from transcription.tex (Danish), 2026-10-13, made in the main session from the collated and
%  page-checked transcription; translation.tex is hand-made, not generated: edit it directly.
%  Conventions: TRANSLATION-PLAYBOOK.md; terms in ebook-method/logik-og-psychologie-1871/GLOSSARY.md (as in the
%  earlier Nielsen translations, Natur og Aand above all: Selvhed = selfhood, Andethed = otherness, Selvmagt =
%  self-power, Magt = power, Anlæg = disposition).  Aand = mind where it is set against the body as a substance or an
%  attribute (Descartes, Spinoza, pp. 1-3), spirit elsewhere (Hans, 2026-10-13).  Page numbers of the print in the
%  margin.  The misprint „Det er er“ (p. 3) is translated as meant.
% =====================================================================
\documentclass[12pt,a4paper,openany]{book}
\usepackage[utf8]{inputenc}
\usepackage[T1]{fontenc}
\usepackage{libertinus}
\usepackage{libertinust1math}
\usepackage{textalpha}
\usepackage[english]{babel}
\usepackage[protrusion=true,expansion=true]{microtype}
\usepackage{setspace}
\usepackage[margin=1.3in]{geometry}
\usepackage{fancyhdr}
\usepackage{hyperref}
\hypersetup{pdftitle={Logic and Psychology. Course for the Academic Year 1871--72}, pdfauthor={Rasmus Nielsen}, hidelinks}
\setlength{\headheight}{15pt}
\pagestyle{fancy}
\fancyhf{}
\fancyhead[LE,RO]{\thepage}
\fancyhead[C]{\textit{R. Nielsen: Logic and Psychology. Course 1871--72}}
\setstretch{1.3}
\usepackage{marginnote}
\newcommand{\opage}[1]{\marginnote{\footnotesize\textit{[#1]}}}

\begin{document}

\begin{titlepage}
\begin{center}
\vspace*{6em}
\rule{3em}{0.4pt}\\[1.5em]
{\LARGE Logic and Psychology.}\\[1.5em]
\rule{3em}{0.4pt}\\[1.5em]
{\large\textbf{Course for the Academic Year 1871--72}}\\[1.5em]
{\small by}\\[1em]
{\Large R. Nielsen.}\\[2em]
\rule{3em}{0.4pt}\\[2em]
\textbf{Copenhagen.}\\[0.5em]
{\small Published by J. H. Schubothe's Bookshop.}\\[3em]
{\small Printed by Louis Klein.}\\[3em]
{\small Translated from the Danish}
\end{center}
\end{titlepage}

% --- p. 1 ---
\opage{1}\begin{center}
{\Large\textbf{Psychology.}}\\[0.4em]
\rule{4em}{0.4pt}\\[0.4em]
Introduction.
\end{center}

What it means in reality that the principle of selfhood must contain the principle for the other, that self-power and
power in the other must confirm each other, that spirit itself must have nature as its ground in order to articulate
and bear the actual natural dispositions, becomes especially evident when we consider

\begin{center}
\textbf{The Relation between Soul and Body.}\\[0.6em]
α) \textbf{Dualism.}
\end{center}

When Cartesius proceeded from the presupposition that soul and body were two heterogeneous substances, the soul the
thinking, the body the extended substance, he thereby set up such a divorce between the essence of selfhood and that
of otherness that interaction became unthinkable. The notion of a mutual physical influence (\textit{systema influxus
physici}) was only a provisional popular assumption, which soon had to give way to metaphysical theories.

In Malebranche an attempt is made to overcome dualism by reducing all causes in things to occasional causes
(\textit{causæ occasionales}) for God as the only acting cause of existence (\textit{causa efficiens}). The relation
between soul and body is then, in consequence, conceived in such a way that God takes occasion from every movement in
the body to move the soul, and conversely (\textit{systema causarum occasionalium}). The almighty Self, by its incessant
intervention in the order of things, gained so absolute an influence over the whole creation that the natural causes
vanished, and with them nature.

% --- p. 2 ---
\opage{2}In order to set a limit to the encroachments of the theory of miracles, Leibniz put forward the doctrine of
pre-established harmony (\textit{Harmonia præstabilita}). The basic thought is that God has from eternity calculated all
the movements in the corporeal and the movements in the soul so exactly that the two otherwise heterogeneous series
correspond to each other point for point.

In Leibniz the principle of selfhood is developed in a peculiar way in the doctrine of monads. All existing substances
are unities of force, each with its disposition to self-development. The soul is thus a monad which, by virtue of its
original disposition, is able from its point of view (\textit{suivant son point de vue}) to represent the whole world
(opposition between microcosm and macrocosm).

But what is matter? In answering this question Leibniz speaks a double language, now taking his stand on the standpoint
of metaphysics, now on that of physics.

In the metaphysical sense matter has no reality; from the subjective side it is only a phenomenon of the obscure,
confused representations of the soul-monad, from the objective side only a complex of the lowest monads
(\textit{Monades nudæ}). From this point of view the whole world is thus a product of nothing but principles of
selfhood.

Considered from the standpoint of physics, on the other hand, matter has an immediate reality; the principle of
otherness comes into its own insofar as Leibniz, in determining the mechanical laws of motion, presupposes pressure,
impact, inertia, etc.

The extreme consequence of dualism is that the two systems of laws, the laws of selfhood (the moral laws of the soul)
and the laws of otherness (the mechanical laws of motion in the body), cannot be reconciled. The notion that the soul
should relate to its body as a master who had an automaton for a servant is so monstrous that it involuntarily demands
the dissolution of dualism.

% --- p. 3 ---
\opage{3}\begin{center}β) \textbf{The Doctrine of Identity.}\end{center}

The foundation is laid by Spinoza. The relative independence that Cartesius had ascribed to the substances, soul and
body, vanishes in relation to God, the absolute substance. As the relative substances lose their independence, they lose
with it their substantiality and are transformed into modifications (\textit{modi}) of one single infinite substance.

It is the one infinite substance, with its attributes corresponding to mind and body, thought and extension, that is
its own cause (\textit{causa sui}) and the cause of all things. One and the same order of things thus constitutes a
series of thinking and a series of extended existences (\textit{res}), the former forming the world of minds, the latter
the world of bodies.

As a mode of the attribute of thought, man is \emph{mind}; as a mode of the attribute of extension, \emph{body}. That by
which man is distinguished from other things is, with respect to extension, a more composite body, and in consequence,
as regards thought, a more composite representation. The mind is the idea of the body (\textit{idea corporis}). If the
human body were a circle, the human mind would be neither more nor less than the idea of this circle.

The doctrine of identity fuses the principles of selfhood and otherness into one single ideal-real principle, and must,
in order to come to terms with the actual oppositions, take refuge in formal circumlocutions and empty verbal
explanations.

If the principle is conceived immediately from the real side, we have materialism. The organic body is a product of the
peculiar interactions of mechanical-chemical causes, and the life of the soul, a modification of these interactions, is
to be explained by a circumlocution of the corporeal (\textit{système de la nature; l’homme machine.}) If the ideal-real
principle is
% --- p. 4 ---
\opage{4}conceived from the ideal point of view, the doctrine of identity becomes speculative.

The sharp dialectical antithesis of \emph{I} and \emph{not-I} (selfhood and otherness), which gives Fichte's Doctrine of
Science its true scientific significance, was in Schelling's philosophy of nature diluted and neutralized to such a
degree that the independence of spirit was lost, and the categories of nature were transformed into hollow symbols in
the doctrine of “\emph{the absolute identity}.”

It is the infinite activity of the All-One that brings forth all life, the natural as well as the ethical, by
accomplishing in an ordered series of stages “\emph{its historical world-process}” and with it its
“\emph{self-liberation}.”

On the first (lowest) stage the principle acts as “\emph{light},” in negative form as “\emph{gravity}”; the product is
“\emph{the inorganic world}.”

On the second stage the principle acts as “\emph{irritable drive},” in negative form as “\emph{sensation},” the product
is “\emph{organic nature}.”

On the third stage the principle acts as “\emph{practical law of reason}” or as “\emph{religion and state},” in negative
form as “\emph{conscious law of drive}”; the product is “\emph{the life of humanity}.”

That such speculative abstractions, however many circumlocutions one employs and however often one varies the worn-out
theme of the all-organizing essence of nature, of the universe that, after many obscure preparations and necessary
dissipations in the elementary, at last gathers itself into organic unity and opens itself up into self-consciousness,
are unfruitful and in themselves unscientific, will be obvious from the present standpoint of science.

\begin{center}γ) \textbf{Identity-Dualism.}\end{center}

What is true in dualism is that soul and body have their origin in opposite, heterogeneous principles, the soul
% --- p. 5 ---
\opage{5}in the principle of selfhood, the body in that of otherness. That matter is selfless in its essence, physics has
settled once and for all in the doctrine of inertia: inertia and selflessness are synonymous categories. It is therefore
just as impossible to make life a resultant of the combined action of physical forces as to derive the self from the
selfless.

What is true in the doctrine of identity is that soul and body, despite their difference in principle, have in the idea
itself their origin in one and the same fundamental principle, \emph{the divine principle of power}.

With full right the representatives of dualism have emphasized that the opposition of soul and body must necessarily
presuppose a divine unity of power; but since the concept of power had not yet been developed in its original
self-duplication, the unity of power was without further ado taken for an immediate expression of the supernatural
creative power; therein consists the fundamental error of dualism.

In the realm of power, in the world of things, the unity of disposition, the power of the idea articulating itself
through acting natural causes, is a unity of opposites of ideal capacity and real subsistence, thus a unity of opposites
of life-principle and matter.

That the life-principle in all its ideality is a principle of power is evident from its acting upon matter: for an
abstract being of thought cannot possibly act upon anything material whatsoever. That matter for its part can be
affected by the life-principle without the substances thereby giving up their physical-chemical nature proves that its
objective truth must be sought in the concept of power which has laid down in the physical properties the possibility
of organic activity.

Dualism rests on the opposition of the duplication of the concept of power through the form of selfhood and that of
otherness; identity, on the other hand, on the reality of the unity of the concept in both forms. The unity of the power
of selfhood in the psychical and the power of otherness in the corporeal does not make the two principles
% --- p. 6 ---
\opage{6}into one principle; but it provides the medium in which the activity of life becomes real, and matter with its
possibilities ideal, the fundamental medium without which an organic interaction between soul and body would be
unthinkable.
(Physiology and psychology).

\begin{center}
\rule{4em}{0.4pt}\\[0.6em]
\textbf{Psychology in Outline.}
\end{center}

It is the life of consciousness with its individual natural ground and organically conditioned functions, whose
transition from sense-consciousness to self-consciousness and spiritual life is the object of psychology. With respect
to this we divide psychology into the doctrine of the natural basic expressions of the self-conscious life, of its
absolute content, and of its absolute validity.

\begin{center}
\textit{a)} \textbf{The Natural Basic Expressions of the Self-Conscious Life:}\\
\textbf{The Faculties of the Soul.}
\end{center}

In sharp contrast to the register-like enumeration of faculties of the soul with which the psychological textbooks are
so often overloaded, some psychologists (Herbart) have flatly denied the concept of a faculty of the soul any reality.
But in principle and disposition the soul is nothing but faculty. The task is to show the different forms and
directions in which the one and the same basic faculty naturally individualizes and ramifies itself.

The most universal form in which selfhood expresses itself is \emph{sensation}: it senses itself in sensing its other,
and it determines its other by sensing itself. Sensation is, in passive immediacy, a sense perception; in sublated
immediacy, a reflection; in active immediacy, a self-determining (dialectic of the active and the passive).

As sense perception, the determinate self-sensation is a \emph{feeling} (the doctrine of the feelings), the sensation
of the other an immediate objectification through the various \emph{senses} (sight, hearing, etc.); as unity of
self-sensation and sensation of the other, it converts its immediacy into
% --- p. 7 ---
\opage{7}\emph{memory} and \emph{imagination}, and herewith reflection begins.

In reflection the sense images are freed by the word (the significance of language) into representations and general
concepts (the life of the understanding) and are clarified into the unity of the subjective and the objective in the
thinking of the universally valid (reason; relation between understanding and reason). Reflection completes itself in
\emph{thinking}. But already the way in which the reciprocal determination between self-sensation and sensation of the
other, between self-determination and determination of the other, is clarified in reflecting and thinking shows that
the passive immediacy of sensation is originally active; that, in other words, self-sensation necessarily presupposes
self-determination.

Originally self-determination, the natural self-expression, is determining as drive; voluntariness proceeds from
involuntary prompting. As the drive points toward the objective, it sets the self in motion; drive is the ground of the
will. But self-determination is reflected and clarifies itself with reflection into will: the will, which is grounded in
the drive, in turn sublates the drive and is grounded in itself (egoism, egoity). The reflected drive is
\emph{inclination}, and the vehement inclination \emph{passion}. The relation between drive and will is seen in the
interplay of passion and will (heroism, asceticism).

\begin{center}
\textit{b)} \textbf{The Absolute Content of the Self-Conscious Life: The Psychology}\\
\textbf{of Ideas.}
\end{center}

Since consciousness is awakened and developed in time by sensuous impressions, finite demands and changing interests,
it is natural that from the beginning it must be filled with finite content. As self-consciousness it nonetheless has an
original self, and precisely for that reason an original claim to a self-valid, absolute content: only the idea, as
idea, can satisfy the hungry subjectivity (spirit and the spiritual).

How naturally the awakening self-consciousness craves
% --- p. 8 ---
\opage{8}the content of the idea is seen from the originality with which every gifted people of nature begins with
religious prefigurations (mythology), ethical fundamental insights, striving to give the supersensible sensuous form and
to determine the universally valid. For the absolute content assumes a fundamentally different form according as it is
originally apprehended by imagination, by thought, or by the will. Form and content so interpenetrate each other that
the one idea, absolute in itself, becomes by means of the three basic forms three basic ideas: beauty, truth, goodness.

\begin{center}
\textit{c)} \textbf{The Absolute Validity of the Self-Conscious Life: The Immortality}\\
\textbf{of the Soul.}
\end{center}

As the organ of eternal ideas, self-consciousness must necessarily have the eternal in its own essence; on this rests
its absolute validity, on this is grounded the doctrine of the immortality of the soul, i.e., of the eternal life of
personality. That the soul is immortal in the metaphysical sense can already be seen from plant and animal life; but the
certainty of an existential and, as regards man, personal immortality can be attained only through clear insight into
the relation between the idea, in particular the idea of the good, and human self-consciousness. Insofar as this
relation in its generality is made clear in science, one may indeed say that the immortality of the soul can be proved
scientifically; but since the conditions of reality for a future life fall outside all experience, and the pure proof
by thought consequently does not reach beyond the metaphysical, the proof is limited, strictly speaking, to establishing
the ethical validity of the demand, not its fulfillment. The perfect certainty of the immortality of the soul is of a
religious-ethical nature and serves precisely to show the difference between faith and knowledge, between practical and
theoretical cognition (the psychologically necessary separation of religion and science).

\end{document}
